% ============================================================================
%  Player-Base Decay Forecasting -- ML Project Checklist Report
%  Framework adapted from A. Geron's "Machine Learning Project Checklist"
%  (Hands-On Machine Learning with Scikit-Learn, Keras & TensorFlow, App. B).
%  Each checklist item records what the project has done so far; items in
%  sections not yet reached are left blank to be filled in later.
%  Compile with pdflatex.
% ============================================================================
\documentclass[11pt,a4paper]{article}

\usepackage[T1]{fontenc}
\usepackage[utf8]{inputenc}
\usepackage{lmodern}
\usepackage{microtype}
\usepackage[margin=2.4cm]{geometry}
\usepackage{amsmath}
\usepackage{xcolor}
\usepackage{enumitem}
\usepackage{booktabs}
\usepackage{parskip}
\usepackage{titlesec}
\usepackage[colorlinks=true]{hyperref}

% ---------------------------------------------------------------------------
% Palette and status badges
% ---------------------------------------------------------------------------
\definecolor{accent}{RGB}{28,60,110}     % deep blue: headings, links
\definecolor{stdone}{RGB}{36,120,66}     % green:  completed
\definecolor{stwip}{RGB}{176,108,12}     % amber:  in progress
\definecolor{sttodo}{RGB}{125,125,125}   % gray:   not started
\definecolor{soft}{RGB}{95,95,95}        % gray text for guidance notes

\hypersetup{linkcolor=accent, urlcolor=accent, citecolor=accent,
            pdftitle={Player-Base Decay Forecasting: ML Project Checklist Report},
            pdfauthor={Sumer Jaitly}}

\titleformat{\section}{\Large\bfseries\color{accent}}{\thesection.}{0.6em}{}
\titlespacing*{\section}{0pt}{1.6em}{0.6em}

\newcommand{\StatusBadge}[2]{\colorbox{#1}{\textcolor{white}{\sffamily\bfseries\footnotesize\;#2\;}}}
\newcommand{\SectionStatus}[2]{\noindent\StatusBadge{#1}{#2}\par\medskip}

% Checklist item: bold prompt, answer follows in normal text
\newlist{checklist}{enumerate}{1}
\setlist[checklist]{label=\textbf{\arabic*.}, ref=\arabic*, leftmargin=2em,
                    itemsep=0.9em, topsep=0.5em}
\newcommand{\chk}[1]{\item \textbf{#1}\par\nopagebreak\vspace{0.15em}}

% Gray italic marker for individual items not yet addressed
\newcommand{\pending}[1][Not yet addressed.]{\textcolor{sttodo}{\itshape #1}}
% Blank writing space for sections not yet reached
\newcommand{\blankanswer}{\vspace{2.8em}}   % TODO: fill in when this step is reached
% Guidance notes carried over from the checklist itself
\newcommand{\guidance}[1]{\textcolor{soft}{\small\itshape Checklist note: #1}\par\medskip}

\setlength{\emergencystretch}{2em}
\raggedbottom

% ============================================================================
\begin{document}

% ---------------------------------------------------------------------------
% Title block
% ---------------------------------------------------------------------------
\begin{center}
  {\color{accent}\rule{\textwidth}{1.2pt}}\\[1.2em]
  {\LARGE\bfseries Player-Base Decay Forecasting for Live-Service Games}\\[0.5em]
  {\large Machine-Learning Project Checklist Report}\\[1.0em]
  Sumer Jaitly \quad$\cdot$\quad \today\\[1.0em]
  {\color{accent}\rule{\textwidth}{1.2pt}}
\end{center}

\vspace{0.5em}
\noindent
This document turns the standard machine-learning project
checklist\footnote{Adapted from Aur\'elien G\'eron, \emph{Hands-On Machine
Learning with Scikit-Learn, Keras \& TensorFlow}, O'Reilly --- Appendix
``Machine Learning Project Checklist''.} into a living report for this
project. Each of the eight steps below reproduces the checklist's items;
under every item we record what the project has actually done. Items in
steps we have not yet reached are left blank, to be filled in as the work
progresses.

\subsection*{Progress at a glance}
\begin{center}
\begin{tabular}{@{}c l l@{}}
\toprule
\textbf{Step} & \textbf{Checklist stage} & \textbf{Status}\\
\midrule
1 & Frame the problem and look at the big picture & \StatusBadge{stdone}{COMPLETED}\\
2 & Get the data                                   & \StatusBadge{stdone}{COMPLETED}\\
3 & Explore the data                               & \StatusBadge{stdone}{COMPLETED}\\
4 & Prepare the data                               & \StatusBadge{stdone}{COMPLETED}\\
5 & Short-list promising models                    & \StatusBadge{stwip}{IN PROGRESS}\\
6 & Fine-tune the system                           & \StatusBadge{sttodo}{NOT STARTED}\\
7 & Present your solution                          & \StatusBadge{stwip}{IN PROGRESS}\\
8 & Launch, monitor, and maintain                  & \StatusBadge{sttodo}{NOT STARTED}\\
\bottomrule
\end{tabular}
\end{center}

\noindent
Steps marked \emph{completed} are complete for the current end-to-end slice
of the project; several are revisited as the modelling iterates (the
checklist is a loop, not a straight line).

% ===========================================================================
\section{Frame the Problem and Look at the Big Picture}
\SectionStatus{stdone}{COMPLETED}

\begin{checklist}

\chk{Define the objective in business terms.}
Build an early-warning system for severe player-base decline in
live-service (``replayable'') games, using Steam as the testbed, so that a
studio, publisher, or live-ops team can prioritise interventions --- content
updates, retention campaigns, pricing changes, community management ---
before a player base becomes unrecoverable. The framing is documented in the
project plan
(\texttt{project\_documents/game\_decay\_forecasting\_project\_plan.qmd})
and in the report draft (\texttt{report/report.qmd}).

\chk{How will your solution be used?}
As decision support for live-ops prioritisation. Because every game has its
own longevity goals, the most flexible primary output is a forecast of the
game's medium-term player-count level (relative to its own history), which
acts as an early-warning signal; the project plan's headline deliverable
builds on this with a calibrated severe-decline risk score and a risk-ranked
watchlist.

\chk{What are the current solutions/workarounds (if any)?}
Manual monitoring of public trackers (SteamDB, SteamCharts) combined with
analyst judgement and ad-hoc threshold rules. These are formalised inside
the project as the baselines any model must beat: a constant
(train-mean) predictor now, with last-value, moving-average, and log-linear
decay baselines planned.

\chk{How should you frame this problem (supervised/unsupervised, online/offline, etc.)?}
Supervised, offline (batch) learning on historical data. The implemented
task is regression on \emph{game-window} examples: a 60-day window of a
game's daily player counts is used to predict the game's average player
count over a 14-day window beginning 60 days after the input window ends,
normalised by the game's \emph{causal peak} (the running maximum of its
14-day rolling average) and modelled in log space:
\[
  y_t \;=\; \log\!\left(
    \frac{\overline{\mathrm{CCU}}^{(14)}_{[\,t+61,\ t+74\,]}}
         {\max_{s\le t}\ \overline{\mathrm{CCU}}^{(14)}_s}
    \;+\; \varepsilon
  \right),
  \qquad \varepsilon = 0.01,
\]
where $\overline{\mathrm{CCU}}^{(14)}$ is the 14-day rolling mean of daily
average concurrent users and $t$ is the last day of the input window. The
project plan frames the eventual headline task as binary severe-decline
classification; the current regression target is its continuous, supporting
variant. Overlapping windows from one game are strongly correlated, so the
effective sample size is roughly the number of games ($\approx$ 2{,}000),
which shapes the validation design throughout.

\chk{How should performance be measured?}
Mean absolute error (MAE) on the log-space target, reported on a held-out
test set against a constant dumb baseline. RMSE was used first and
deliberately abandoned: the target is extremely heavy-tailed (values up to
$60\times$ the causal peak), and under squared error the worst
$\sim$1\% of points contributed $\sim$86\% of the total, swamping genuine
skill. Log-space MAE reads directly as a typical multiplicative error. For
the planned classifier stage: PR-AUC, precision@top-$K$, and calibration.

\chk{Is the performance measure aligned with the business objective?}
Yes. Log-space MAE is scale-free, so mega-hits and niche games weigh
comparably, and it answers ``how far off, proportionally, is the forecast''
--- the quantity a per-game intervention decision needs. Precision@top-$K$
(planned) matches a live-ops team's limited intervention budget.

\chk{What would be the minimum performance needed to reach the business objective?}
The qualitative bar is fixed by the project plan: the model must clearly
beat the naive baselines and simple threshold rules to justify existing.
\pending[A quantitative minimum (e.g.\ a target multiplicative error) has
not yet been set.]

\chk{What are comparable problems? Can you reuse experience or tools?}
Customer-churn prediction, demand forecasting, and survival analysis;
technically, sliding-window time-series regression. Reused: the standard
Python stack (pandas, NumPy, scikit-learn, matplotlib), this checklist's
workflow, and the dataset's companion research (Vihanga~et~al., IEEE CoG
2019, DOI \texttt{10.1109/CIG.2019.8848108}) on player-population patterns.

\chk{Is human expertise available?}
The author's own domain familiarity with Steam and live-service player-base
dynamics (used, e.g., in the \emph{Among Us} case study that stress-tested
the definition of ``collapse''). No external domain expert is engaged.

\chk{How would you solve the problem manually?}
As an analyst would: plot the trajectory, compare the current level with the
historical peak, estimate the recent trend and volatility, and extrapolate.
This manual recipe directly motivated the five window-summary features
(mean, standard deviation, max, min, net change).

\chk{List the assumptions you or others have made so far.}
(i)~Pre-COVID dynamics are representative --- data are clipped at
2020-03-14; (ii)~concurrent players (CCU) proxies player-base health (it is
\emph{not} daily active users); (iii)~the sample --- the top 2{,}000 games by
player count on 11~Dec~2017 --- carries survivorship/selection bias, so
conclusions may not transfer to smaller games; (iv)~the question is
meaningful only for replayable games; (v)~windows from the same game are
strongly correlated; (vi)~causal-peak normalisation makes games of different
scales comparable; (vii)~a 60-day-ahead forecast is early enough to act on.

\chk{Verify assumptions if possible.}
Verified where the data allow: the COVID player-count inflation is visible
in the series (hence the clip); the heavy tail of the target was confirmed
by inspection of its distribution (motivating the log target and MAE); the
same-game-correlation assumption is enforced structurally by a game-level
train/test split; the dataset's two granularity tiers (5-minute vs hourly)
are equalised by daily aggregation before any modelling.

\end{checklist}

% ===========================================================================
\section{Get the Data}
\SectionStatus{stdone}{COMPLETED}
\guidance{automate as much as possible so you can easily get fresh data.}

\begin{checklist}

\chk{List the data you need and how much you need.}
Multi-year, daily-resolution player-count histories for a large and varied
set of games (thousands, since the effective sample size is the number of
games), plus per-game metadata (tags, genres, prices) for later feature
work.

\chk{Find and document where you can get that data.}
Mendeley Data: \emph{``Steam Games Dataset: Player count history, Price
history and data about games''}, Wannigamage, Barlow, Lakshika \& Kasmarik
(UNSW Canberra), 2020, DOI \texttt{10.17632/ycy3sy3vj2.1}, CC~BY~4.0
(\url{https://data.mendeley.com/datasets/ycy3sy3vj2/1}). Chosen after a
documented survey of alternatives (monthly SteamCharts derivatives, snapshot
datasets, live trackers without bulk export), none of which offer daily,
multi-game history.

\chk{Check how much space it will take.}
Raw data $\approx$ 13\,GB unpacked (Part~1: 5-minute samples for the top
1{,}000 games; Part~2: hourly for the next 1{,}000). Processed daily
aggregates: 54\,MB.

\chk{Check legal obligations, and get the authorization if necessary.}
Licensed CC~BY~4.0; attribution is given in the notebook, the report, and
this document. No further authorisation required.

\chk{Get access authorizations.}
None needed --- public dataset.

\chk{Create a workspace (with enough storage space).}
Local git repository \texttt{player\_forecasting/} with
\texttt{data/mendeley/} (raw, kept out of version control) and
\texttt{data/avg\_daily/} (processed), plus \texttt{notebooks/},
\texttt{report/}, and \texttt{project\_documents/}.

\chk{Get the data.}
Downloaded and unpacked: \texttt{PlayerCountHistory} Parts 1--2,
application metadata (information, genres, tags, developers, publishers),
and price history.

\chk{Convert the data to a format you can easily manipulate (without changing the data itself).}
Per-game 5-minute/hourly CSVs are aggregated to per-game \emph{daily
average} CSVs in \texttt{data/avg\_daily/} (one file per game), then
combined into a single games\,$\times$\,days matrix for modelling. Raw
files are kept intact.

\chk{Ensure sensitive information is deleted or protected.}
The data are aggregate, public concurrent-player counts; no personal or
sensitive information is present.

\chk{Check the size and type of data.}
Time series: 2{,}000 games $\times$ $\sim$970 days (2017-12-14 to
2020-08-12), two sampling tiers equalised by daily aggregation. Analysis is
restricted to the 822 pre-COVID days (up to 2020-03-14). Counts are
heavy-tailed, spanning several orders of magnitude across games.

\chk{Sample a test set, put it aside, and never look at it.}
10\% of games (197 of the 1{,}997 usable games) are held out via a seeded
game-level split \emph{before} any modelling. Splitting by game --- rather
than by window --- guarantees that overlapping windows from one game never
straddle the train/test boundary (no leakage). The test set is consulted
only to report final numbers.

\end{checklist}

% ===========================================================================
\section{Explore the Data}
\SectionStatus{stdone}{COMPLETED}
\guidance{try to get insights from a field expert for these steps.}

\begin{checklist}

\chk{Create a copy of the data for exploration.}
All exploration works on derived copies (daily aggregates and in-memory
dataframes); the raw Mendeley files are never modified.

\chk{Create a Jupyter notebook to keep record of your data exploration.}
\texttt{notebooks/data\_engineering.ipynb} records the full exploration and
pipeline narrative; \texttt{data\_engineering.py} has been started to hold
the reusable pieces.

\chk{Study each attribute and its characteristics.}
\begin{itemize}[leftmargin=1.4em, itemsep=0.2em]
  \item \emph{Playercount} --- float concurrent users, $\ge 0$, unbounded
        above; heavy-tailed (roughly logarithmic) across games; missing
        entries at 5-minute granularity (dropped before aggregation) and
        occasional full-day gaps (become NaNs, handled at windowing);
        strong day-to-day jitter, studied via smoothing comparisons
        (exponential moving average vs rolling mean).
  \item \emph{Time} --- 5-minute/hourly timestamps, normalised to dates.
  \item \emph{Tags} --- multi-valued categorical strings per game;
        lower-cased, frequency-counted, top-30 inspected as candidate
        features.
\end{itemize}

\chk{For supervised learning tasks, identify the target attribute(s).}
Identified and implemented: the log-space, causally-normalised 14-day future
average defined in \S1, item~4.

\chk{Visualize the data.}
Per-game trajectory plots; raw vs EMA vs rolling-average smoothing
comparisons; histograms of the five engineered features; histogram of the
target distribution (which exposed the heavy tail).

\chk{Study the correlations between attributes.}
Feature distributions have been studied individually.
\pending[A systematic correlation / feature-redundancy analysis is pending.]

\chk{Study how you would solve the problem manually.}
Done --- see \S1, item~10; the manual recipe seeded the feature set.

\chk{Identify the promising transformations you may want to apply.}
Daily aggregation; 14-day rolling mean (jitter removal); causal-peak
(expanding-maximum) normalisation; log transform of the target; log
transform of the features (validated in supplementary benchmarks --- see
\S5).

\chk{Identify extra data that would be useful.}
Game tags (frequency analysis done; one-hot encoding planned as the next
feature source); price history (available in the dataset, leak-safe from
2019-04); Twitch viewership noted in the report as a possible extension.

\chk{Document what you have learned.}
Documented across the notebook narrative, the Quarto report draft
(\texttt{report/report.qmd}), the project plan, and this checklist report.

\end{checklist}

% ===========================================================================
\section{Prepare the Data}
\SectionStatus{stdone}{COMPLETED}
\guidance{work on copies of the data; write functions for all
transformations so they can be reapplied to fresh data, future projects,
the test set, and new instances, and so preparation choices can be treated
as hyperparameters.}

All transformations are implemented as reusable functions
(\texttt{to\_daily\_avg}, the window \texttt{extractor\_log},
\texttt{to\_feature}); migration of these from the notebook into an
importable module (\texttt{data\_engineering.py}) has begun.

\begin{checklist}

\chk{Data cleaning.}
Rows with null player counts are dropped before aggregation; three games
with fewer than 50 clean windows are removed (1{,}997 games remain); all
series are clipped at 2020-03-14 to exclude the COVID regime; the window
extractor skips any window contaminated by NaNs (sliding 14 days at a time,
jumping 60 days when the input window itself is contaminated).

\chk{Feature selection (optional).}
The current slice deliberately uses player-count history only. Degenerate
windows (window mean, max, or min exactly zero --- games effectively dead or
unobserved) are filtered out of the design matrices.

\chk{Feature engineering, where appropriate.}
Daily average CCU $\rightarrow$ 14-day rolling mean $\rightarrow$
causal-peak normalisation $\rightarrow$ 60-day input windows
$\rightarrow$ five summary features per window: mean, standard deviation,
maximum, minimum, and net change (last $-$ first). The target is
log-transformed with $\varepsilon = 0.01$ (\S1, item~4). Log-transforming
the four positive \emph{features} as well has been validated in
supplementary experiments (\S5, item~3) and is queued for adoption in the
notebook.

\chk{Feature scaling: standardize or normalize features.}
Causal-peak normalisation already places all games on a comparable
$[0,1]$-ish scale, and the target is log-scaled. Explicit standardisation is
applied where the algorithm requires it (kNN, SVR, and kernel-approximation
models in the supplementary benchmarks); tree and linear models run on the
features directly.

\end{checklist}

% ===========================================================================
\section{Short-List Promising Models}
\SectionStatus{stwip}{IN PROGRESS}
\guidance{if the data is huge, sample smaller training sets; automate as
much as possible.}

\begin{checklist}

\chk{Train many quick and dirty models from different categories.}
In the notebook: a constant (train-mean) dumb baseline and ordinary linear
regression on the five features ($\sim$89k train / $\sim$9.7k test window
examples). In supplementary benchmark scripts run on the identical split:
ridge regression on the raw 60-day window, $k$-nearest neighbours
($k{=}20$), an RBF-kernel approximation + ridge, an RBF-kernel SVR (on an
8k subsample --- full kernel methods do not scale to 89k points), and
histogram gradient-boosted trees (on raw and log targets, raw and log
features).

\chk{Measure and compare their performance.}
Held-out test MAE in log space: dumb baseline 0.933 (0.899 for the
MAE-optimal constant, the train median); linear regression 0.530; gradient
boosting 0.304; linear regression on log-transformed features 0.306.
Interpreted multiplicatively: typical forecast error falls from
$\times 2.5$ (dumb) to $\times 1.7$ (linear) to $\times 1.36$ (gradient
boosting, or linear-on-log-features).
\pending[$N$-fold cross-validation with mean~$\pm$~std per model is
pending; all figures so far come from a single fixed game-level split.]

\chk{Analyze the most significant variables for each algorithm.}
The dominant structure identified so far is \emph{log-log linearity}:
log-transforming the positive features moves linear regression from 0.530
to 0.306, statistically matching gradient boosting --- i.e.\ the
nonlinearity the trees exploit is essentially logarithmic curvature in the
feature--target relationship. \pending[Formal per-model feature-importance
analysis is pending.]

\chk{Analyze the types of errors the models make.}
Under squared error, unpredictable ``breakout'' games (future counts far
above the historical peak) dominated: the worst $\sim$1\% of test points
carried $\sim$86\% of the squared error, which motivated the switch to the
log target and MAE. A human analyst would avoid these errors only with
information from outside the count history --- patch and content-update
schedules, sales, streamer coverage --- which is exactly the direction of
the planned tag/price/exogenous features.

\chk{Have a quick round of feature selection and engineering.}
Completed once: raw daily inputs were replaced by 14-day rolling,
causally-normalised inputs (removing jitter), the raw ratio target was
replaced by its log, and log features were identified as the next
adoption.

\chk{Have one or two more quick iterations of the five previous steps.}
One full loop is complete (raw target/features $\rightarrow$
smoothed/log). \pending[Next loop planned: add tag features and re-run the
comparison.]

\chk{Short-list the top three to five most promising models.}
Current short list: (i)~\textbf{linear regression on log features} ---
interpretable and within noise of the best; (ii)~\textbf{histogram
gradient-boosted trees} --- equally strong, handles NaNs and categorical
(tag) features natively, the natural vehicle for the next feature round.
Kernel SVMs were evaluated and ruled out (no accuracy gain over linear at
$\mathcal{O}(n^2)$ kernel cost); $k$NN underperformed both finalists.

\end{checklist}

% ===========================================================================
\section{Fine-Tune the System}
\SectionStatus{sttodo}{NOT STARTED}
\guidance{use as much data as possible; automate what you can.}

\begin{checklist}

\chk{Fine-tune the hyperparameters using cross-validation.}
% Treat data-transformation choices as hyperparameters; prefer random
% search over grid search (or Bayesian optimisation if training is slow).
\blankanswer

\chk{Try ensemble methods.}
% Combining the best models will often perform better than running them
% individually.
\blankanswer

\chk{Measure the final model's performance on the test set to estimate the generalization error.}
\blankanswer

\end{checklist}

\begin{quote}\itshape\color{soft}
Don't tweak your model after measuring the generalization error: you would
just start overfitting the test set.
\end{quote}

% ===========================================================================
\section{Present Your Solution}
\SectionStatus{stwip}{IN PROGRESS}

\begin{checklist}

\chk{Document what you have done.}
Ongoing throughout: the project plan (Quarto, with HTML/PDF renders), the
technical report draft (\texttt{report/report.qmd} --- objective, scope,
definition of ``player-base collapse'', dataset description, target
definition, processing steps, with bibliography), the narrated pipeline
notebook, and this checklist report.

\chk{Create a nice presentation (highlight the big picture first).}
\blankanswer

\chk{Explain why your solution achieves the business objective.}
Argued in draft form in the report: a continuous forecast of future
player-count level (relative to a game's own peak) is the most flexible
early-warning input because longevity goals differ per game.
\pending[To be completed once final results are in.]

\chk{Present interesting points you noticed along the way.}
Collected so far --- \emph{what worked:} daily aggregation, 14-day rolling
smoothing, causal-peak normalisation, the log target, and the game-level
split. \emph{What did not:} RMSE as the headline metric (outlier-dominated),
kernel SVMs (no gain over linear, poor scaling), and inverting a
raw-feature linear model trained in log space back to the original scale.
Assumptions and limitations are listed in \S1 and in the report's
Limitations section.

\chk{Ensure key findings are communicated through beautiful visualizations or easy-to-remember statements.}
Candidate headline statements are drafted (``typical multiplicative
forecast error $\times 1.36$, versus $\times 2.5$ for the dumb baseline'';
``one log transform is worth a gradient-boosted model'').
\pending[Final presentation-quality visualisations are pending.]

\end{checklist}

% ===========================================================================
\section{Launch, Monitor, and Maintain}
\SectionStatus{sttodo}{NOT STARTED}

\begin{checklist}

\chk{Get your solution ready for production.}
% Plug into production data inputs, write unit tests, etc.
\blankanswer

\chk{Write monitoring code to check live performance and trigger alerts.}
% Beware slow degradation (model rot); monitor input quality too.
\blankanswer

\chk{Retrain your models on a regular basis on fresh data.}
% Automate as much as possible.
\blankanswer

\end{checklist}

\end{document}
